\section{What testing and search can establish}
\label{sec:search}

\subsection{A counterexample and an absence are asymmetric}
A verified bypass refutes the associated boundary claim. Failure to find one does not prove completeness in an open world. Any finite collection of observed routes is compatible with a larger world containing an additional unobserved route, unless the environment's structure independently excludes it. Complete finite-state exploration can close this gap \emph{within its model}; ordinary red-teaming cannot claim that same scope.

The reporting process introduces another ambiguity. Consider two worlds with identical reported cases. In the first, the reports exhaust the relevant family; in the second, additional cases exist but the searcher never reports them. Every statistic computed only from the reports has the same distribution in the two worlds. It cannot distinguish exhaustion from systematic omission. The omission may be deliberate, a shared blind spot, or lack of search capability. An honest searcher is therefore insufficient unless it also has justified coverage and reports on the relevant distribution.

Species-discovery estimates and capture--recapture methods can be useful diagnostics under their sampling assumptions. They do not remove this identifiability problem. More independent-looking reporters are not a certificate that one has covered a withheld family. Likewise, a more capable search model does not automatically dominate a deployed model on every rare case.

\subsection{A defensible finite-sample statement}
Fix the evaluated protocol, opponent generator, task distribution, and number of trials before evaluating the final holdout. If the episodes are independent Bernoulli trials with failure probability $p$, observing zero failures in $n$ trials gives the exact one-sided $(1-\delta)$ confidence upper limit \cite{clopper}
\begin{equation}
p_U=1-\delta^{1/n}.
\label{eq:binomial}
\end{equation}
At $n=3{,}000$ and $\delta=.05$, $p_U\simeq9.98\times10^{-4}$. To obtain $p_U\le10^{-6}$ after zero failures requires at least $2{,}995{,}731$ such trials. These are confidence limits under the sampling model, not posterior probabilities that a system is safe.

The unit matters. Ten thousand actions from one coordinated campaign are not ten thousand independent campaigns. If the target is episode-level joint failure, the experimental unit must preserve coordination, shared state, and attacker adaptation within the episode. Adaptive attack discovery belongs in development; the final held-out protocol must be defined separately. Repeatedly checking the ordinary interval and stopping when it looks favorable invalidates its nominal coverage; use a fixed design or a justified sequential method.

\subsection{Transfer is another assumption}
Let $Q$ be the evaluation distribution over complete episodes and $P$ the deployment distribution. For any failure event $F$, a justified total-variation bound $\operatorname{TV}(P,Q)\le d$ implies
\[
P(F)\le Q(F)+d.
\]
Alternatively, $P\ll Q$ and a density-ratio bound $dP/dQ\le\alpha$ imply $P(F)\le\alpha Q(F)$. These are elementary change-of-measure statements. The important words are \emph{justified} and \emph{complete episodes}. The data sampled from $Q$ do not alone bound $d$ or $\alpha$, particularly for unseen inputs or stronger adaptive opponents.

Planted failures can validate instrumentation and estimate recall for the planted family. They do not, without a transfer argument, establish recall for unknown failure families. Likewise, a Bayesian calculation about silent search requires an explicit prior and a justified likelihood for silence when a bypass exists. Neither an honest searcher nor a high capability score supplies those assumptions automatically.

Our reporting rule is consequently strict: state the tested scope and its interval; state any transfer assumption separately; do not manufacture an open-world risk estimate when the transfer is unknown. Testing remains worthwhile because it finds counterexamples, compares interventions, and supports bounded deployment decisions. It simply cannot carry a claim larger than its evidence.
